\subsection*{B.1\quad Intellectual property}
\emph{Core value-capture transition: from creators of content to owners of IP.}

Policies, programmes, laws, institutions or practical services that help creative SMEs identify, document, protect, retain, exploit and manage intangible creative assets so they can earn income beyond a one-off sale. This includes copyright, related rights, trademarks, design rights, patents, geographical indications, trade secrets, contracts, chain of title, licensing, royalties, collective management, anti-piracy, enforcement, rights metadata and IP-backed or royalty-based monetisation.

\textbf{Terms that signal the channel:} copyright; intellectual property; authors' rights; related rights; neighbouring rights; trademark; design right; patent; geographical indication; trade secret; licensing; royalties; royalty collection; collective management organization; CMO; anti-piracy; enforcement; rights retention; contract; chain of title; ownership; intangible asset; IP-backed lending; rights metadata; registry; creator compensation; fair remuneration.

\textbf{Excluded or down-weighted:} generic cultural heritage, creativity, innovation or entrepreneurship language without ownership, protection, licensing, royalties, contracts or monetisation of intangible assets.

\subsection*{B.2\quad Brand development and positioning}
\emph{Core value-capture transition: from suppliers to builders of brands.}

Policies or services that help creative SMEs convert distinctiveness, origin, quality, identity and reputation into market recognition, price premia, repeat demand and direct audience or customer relationships. This includes brand strategy, positioning, packaging, storytelling, marketing, social media presence, owned channels, trademarks when used for brand identity, reputation-building, buyer recognition, diaspora/audience development and visible crediting of creators or firms.

\textbf{Terms that signal the channel:} brand; branding; brand identity; positioning; market recognition; reputation; origin story; storytelling; packaging; label; logo; trademark for brand identity; price premium; willingness to pay; repeat demand; customer loyalty; audience relationship; owned channel; mailing list; social media presence; visible credit; promotion; influencer; diaspora marketing; cultural distinctiveness.

\textbf{Excluded or down-weighted:} generic SME support, start-up support, incubators, productivity or scaling language without recognition, positioning, reputation, customer relationship or price-premium logic.

\subsection*{B.3\quad Market leverage and access}
\emph{Core value-capture transition: from intermediary dependence to market leverage.}

Policies, infrastructure or services that help creative SMEs reach buyers and audiences on better terms, diversify routes to market, reduce dependence on gatekeepers and intermediaries, negotiate deliberately, control customer access or data, and build bargaining power individually or collectively. This includes export promotion, showcases, fairs, mobility, distribution, e-commerce routes, direct-to-consumer sales, platform access, buyer linkages, logistics, customs, certification for market entry and collective bargaining or producer organizations.

\textbf{Terms that signal the channel:} market access; export promotion; distribution; route to market; buyer linkage; audience access; fair; showcase; festival; mobility; touring; residency; direct-to-consumer; direct-to-fan; platform access; marketplace; intermediary; gatekeeper; bargaining; negotiation; collective bargaining; producer organization; cooperative; customer data; first-party data; distributor; sales agent; customs; logistics; parcel shipment; market intelligence.

\textbf{Excluded or down-weighted:} generic trade openness or infrastructure language that does not affect creative SMEs' access to buyers, audiences, distribution channels or negotiating position.

\subsection*{B.4\quad Digital capabilities and commerce}
\emph{Enabling condition: digital tools that make ownership, branding and market leverage cheaper or more feasible.}

Policies or programmes that improve creative SMEs' capacity to use digital tools for creation, rights management, audience-building, commerce, analytics and distribution. The definition should privilege firm-facing digital adoption that helps SMEs own assets, build brands, reach customers, collect revenues, monitor use and reduce intermediary dependence.

\textbf{Terms that signal the channel:} digital tool; digital transformation; e-commerce; online sales; platform; website; social media; newsletter; community platform; streaming; direct-to-fan; digital rights management; DRM; content ID; online registry; metadata; licensing marketplace; royalty tracking; usage monitoring; analytics; first-party data; search engine optimization; online marketing; AI translation; subtitling; dubbing; digital payment; broadband; internet.

\textbf{Excluded or down-weighted:} generic government digitalization, statistics systems or public-sector ICT modernization without a firm-facing link to creative production, IP, branding, commerce, payments, distribution or analytics.

\subsection*{B.5\quad Finance and financial capability}
\emph{Enabling condition: the right money at the right stage for intangible-rich creative SMEs.}

Policies, funds, instruments or advisory services that help creative SMEs finance the creation, protection, monetisation and scaling of intangible assets without losing ownership or strategic control. This includes grants for rights registration, concessional finance, working capital, invoice factoring, trade finance, export credit, credit guarantees, crowdfunding, royalty/revenue-based finance, IP-backed lending, investment readiness, business planning, financial records and separating business from personal accounts.

\textbf{Terms that signal the channel:} access to finance; grant; concessional finance; loan; credit; line of credit; guarantee; equity; investment readiness; business plan; financial records; working capital; invoice factoring; receivables; trade finance; export credit; credit insurance; crowdfunding; revenue-based finance; royalty financing; IP-backed lending; collateral; cash flow; payment delay; bank account; savings account; financial management.

\textbf{Excluded or down-weighted:} generic public spending, subsidies or macro-finance language that does not improve creative SMEs' access to capital, cash-flow management, investment readiness or intangible-asset monetisation.

\subsection*{B.6\quad Sustainability, standards and certification}
\emph{Enabling condition: sustainability as credible differentiation and market access, not only compliance.}

Policies, standards, certifications or services that help creative SMEs adopt verifiable environmental, social, quality or safety practices, including eco-labelling, traceability, sustainable materials, circular production, energy management, clean energy, audits, certification, compliance systems and buyer-trust standards. The emphasis is on concrete firm practices and credible standards, not broad sustainable-development rhetoric.

\textbf{Terms that signal the channel:} sustainability; environmental practice; social standard; certification; quality certification; eco-label; traceability; responsible sourcing; circular design; circular economy; sustainable materials; waste management; recycling; carbon; emissions; energy management; clean energy; solar; green production; compliance; due diligence; audit; inspection; product passport; fair trade; buyer trust; anti-greenwashing.

\textbf{Excluded or down-weighted:} generic sustainable development, cultural diversity, social inclusion or environmental rhetoric without concrete firm practices, energy management, clean energy, standards, certification, compliance, traceability, audits or buyer credibility.

